% from ega3-1-fr.tex
\oldpage[0\textsubscript{III}]{70}
On a donc par hypothèse
$\mathrm F^{k_0}(A_k)=A_k$ et $\mathrm F^{k+1}(A_k)=0$, autrement dit la
filtration considérée est \emph{finie}
(\hyperref[0.11.1.3-fr]{11.1.3}). Les objets gradués associés à cette
filtration sont donc
\[
\operatorname{gr}^p(A_k)=
\operatorname{Ker}(A_k\longrightarrow A_{p-1})
/\operatorname{Ker}(A_k\longrightarrow A_p)
\]
et par suite, $\operatorname{gr}^p(A_k)$ est isomorphe à l'image par
$A_k\to A_p$ de $\operatorname{Ker}(A_k\to A_{p-1})$~; en vertu de la
transitivité des morphismes définissant un système projectif, on a donc
\begin{equation}
\label{0.13.4.1.2-fr}
\tag{13.4.1.2}
\operatorname{gr}^p(A_k)=
\operatorname{Ker}(A_p\longrightarrow A_{p-1})
\cap\operatorname{Im}(A_k\longrightarrow A_p)
\end{equation}
mais comme, en vertu de
(\hyperref[0.13.4.1.1-fr]{13.4.1.1}), on a
$\operatorname{Ker}(A_p\to A_{p-1})=\operatorname{gr}^p(A_p)$, on a aussi
\begin{equation}
\label{0.13.4.1.3-fr}
\tag{13.4.1.3}
\operatorname{gr}^p(A_k)=
\operatorname{gr}^p(A_p)\cap\operatorname{Im}(A_k\longrightarrow A_p).
\end{equation}
Les définitions précédentes montrent en outre que l'on a pour $k\leq h$
\[
u_{kh}(\mathrm F^p(A_h))\subset\mathrm F^p(A_k)
\]
et par suite que les $\operatorname{gr}^p(u_{kh})$ définissent un
\emph{système projectif}
$(\operatorname{gr}^p(A_k))_{k\in\mathbb Z}$ pour tout $p\in\mathbb Z$.
\end{env}

\begin{env}[13.4.2]
\label{0.13.4.2-fr}
Nous dirons que le système projectif $\mathbf A$ est
\emph{essentiellement constant} si les morphismes
$A_{k+1}\to A_k$ sont des \emph{isomorphismes} pour $k$ assez grand. Nous
dirons que le système projectif $\mathbf A$ est \emph{strict} si les
morphismes $A_i\to A_j$ $(j\leq i)$ sont des \emph{épimorphismes}.
Lorsque $\mathbf A$ est strict, il résulte de
(\hyperref[0.13.4.1.3-fr]{13.4.1.3}) que pour $p\leq k\leq h$, le morphisme
canonique $\operatorname{gr}^p(A_h)\to\operatorname{gr}^p(A_k)$ est un
\emph{isomorphisme}, autrement dit, le système projectif
$(\operatorname{gr}^p(A_k))_{k\in\mathbb Z}$ est essentiellement constant.
La suite des objets $\operatorname{gr}^p(A_p)$ (identifiés à
$\varprojlim_k\operatorname{gr}^p(A_k)$ pour tout $p$) se note alors
$\operatorname{gr}^\bullet(\mathbf A)$ et s'appelle \emph{l'objet gradué
associé au système projectif strict} $\mathbf A=(A_k)$.

Si l'on suppose maintenant que le système projectif $\mathbf A$ (limité
inférieurement) vérifie (ML), on sait
(\hyperref[0.13.1.2-fr]{13.1.2}) que le système projectif
$\mathbf A'=(A'_k)$ des objets d'«~images universelles~» est \emph{strict},
et il est par ailleurs limité inférieurement~; l'objet gradué
$\operatorname{gr}^\bullet(\mathbf A')$ associé à $\mathbf A'$ est alors
encore appelé l'objet gradué \emph{associé à} $\mathbf A$ et noté
$\operatorname{gr}^\bullet(\mathbf A)$.
\end{env}

\begin{proposition}[13.4.3]
\label{0.13.4.3-fr}
\emph{Soit $\mathbf A=(A_k)_{k\in\mathbb Z}$ un système projectif limité
inférieurement dans une catégorie abélienne $C$. Les deux conditions
suivantes sont équivalentes :}
\begin{enumerate}
\item[\emph{a)}] \emph{$\mathbf A$ vérifie la condition (ML).}
\item[\emph{b)}] \emph{Pour tout $p\in\mathbb Z$, le système projectif
$(\operatorname{gr}^p(A_k))_{k\in\mathbb Z}$ est essentiellement constant.}
\end{enumerate}
\emph{En outre, lorsque ces conditions sont satisfaites, on a pour tout
$p\in\mathbb Z$ un isomorphisme canonique}
\begin{equation}
\label{0.13.4.3.1-fr}
\tag{13.4.3.1}
\operatorname{gr}^p(\mathbf A)\xrightarrow{\sim}
\varprojlim_k\operatorname{gr}^p(A_k).
\end{equation}
\end{proposition}

\begin{proof}
Il résulte aussitôt de
(\hyperref[0.13.4.1.2-fr]{13.4.1.2}) que \emph{a)} implique \emph{b)}~; la
même formule appliquée au système projectif $\mathbf A'$ (notations de
(\hyperref[0.13.4.2-fr]{13.4.2})) donne l'isomorphisme
(\hyperref[0.13.4.3.1-fr]{13.4.3.1}) par définition. Pour $k\leq h$,
posons $A_{kh}=\operatorname{Im}(A_h\to A_k)$~; si $k\leq h\leq j$, on a
$A_{kj}\subset A_{kh}\subset A_k$. Munissons $A_{kh}$ de la filtration
induite par $(\mathrm F^p(A_k))$~; on vérifie aussitôt, en vertu de la
transitivité des morphismes définissant $\mathbf A$, que cette filtration
est aussi la filtration quotient de $(\mathrm F^p(A_h))$~; par suite, on a
\begin{equation}
\label{0.13.4.3.2-fr}
\tag{13.4.3.2}
\operatorname{gr}^p(A_{kh})=
\operatorname{Im}\bigl(\operatorname{gr}^p(A_h)\longrightarrow
\operatorname{gr}^p(A_k)\bigr).
\end{equation}
